\documentclass[10pt,a4paper]{amsart}
\usepackage{amsmath,amssymb,amsthm,graphicx}
\usepackage{verbatim}
\newenvironment{reference}{\comment}{\endcomment}
%\newenvironment{reference}{}{}
\newenvironment{slogan}{\comment}{\endcomment}
\newenvironment{history}{\comment}{\endcomment}

% For commutative diagrams we use Xy-pic
\usepackage[all]{xy}

% We use 2cell for 2-commutative diagrams.
\xyoption{2cell}
% Theorem environments.
%
\theoremstyle{plain}
\newtheorem{theorem}[subsection]{Theorem}
\newtheorem{proposition}[subsection]{Proposition}
\newtheorem{lemma}[subsection]{Lemma}

\theoremstyle{definition}
\newtheorem{definition}[subsection]{Definition}
\newtheorem{example}[subsection]{Example}
\newtheorem{exercise}[subsection]{Exercise}
\newtheorem{situation}[subsection]{Situation}

\theoremstyle{remark}
\newtheorem{remark}[subsection]{Remark}
\newtheorem{remarks}[subsection]{Remarks}

\numberwithin{equation}{subsection}

% Macros
%
\def\lim{\mathop{\mathrm{lim}}\nolimits}
\def\colim{\mathop{\mathrm{colim}}\nolimits}
\def\Spec{\mathop{\mathrm{Spec}}}
\def\Hom{\mathop{\mathrm{Hom}}\nolimits}
\def\Ext{\mathop{\mathrm{Ext}}\nolimits}
\def\SheafHom{\mathop{\mathcal{H}\!\mathit{om}}\nolimits}
\def\SheafExt{\mathop{\mathcal{E}\!\mathit{xt}}\nolimits}
\def\Sch{\mathit{Sch}}
\def\Mor{\mathop{\mathrm{Mor}}\nolimits}
\def\Ob{\mathop{\mathrm{Ob}}\nolimits}
\def\Sh{\mathop{\mathit{Sh}}\nolimits}
\def\NL{\mathop{N\!L}\nolimits}
\def\CH{\mathop{\mathrm{CH}}\nolimits}
\def\proetale{{pro\text{-}\acute{e}tale}}
\def\etale{{\acute{e}tale}}
\def\QCoh{\mathit{QCoh}}
\def\Ker{\mathop{\mathrm{Ker}}}
\def\Im{\mathop{\mathrm{Im}}}
\def\Coker{\mathop{\mathrm{Coker}}}
\def\Coim{\mathop{\mathrm{Coim}}}

% Boxtimes
%
\DeclareMathSymbol{\boxtimes}{\mathbin}{AMSa}{"02}

%
% Macros for moduli stacks/spaces
%
\def\QCohstack{\mathcal{QC}\!\mathit{oh}}
\def\Cohstack{\mathcal{C}\!\mathit{oh}}
\def\Spacesstack{\mathcal{S}\!\mathit{paces}}
\def\Quotfunctor{\mathrm{Quot}}
\def\Hilbfunctor{\mathrm{Hilb}}
\def\Curvesstack{\mathcal{C}\!\mathit{urves}}
\def\Polarizedstack{\mathcal{P}\!\mathit{olarized}}
\def\Complexesstack{\mathcal{C}\!\mathit{omplexes}}
% \Pic is the operator that assigns to X its picard group, usage \Pic(X)
% \Picardstack_{X/B} denotes the Picard stack of X over B
% \Picardfunctor_{X/B} denotes the Picard functor of X over B
\def\Pic{\mathop{\mathrm{Pic}}\nolimits}
\def\Picardstack{\mathcal{P}\!\mathit{ic}}
\def\Picardfunctor{\mathrm{Pic}}
\def\Deformationcategory{\mathcal{D}\!\mathit{ef}}
\usepackage{cleveref}
\crefname{originallemma}{original lemma}{original lemmas}
\Crefname{originallemma}{Original lemma}{Original lemmas}
\makeatletter
\newlabel{lemma-finite-presentation-local}{{0410}{source}{}{sourcepointer-1}{}}
\newlabel{lemma-finite-presentation-local@cref}{{[originallemma][0][]0410}{[1][0][]source}{}{}{}}
\makeatother
\crefname{originallemma}{original lemma}{original lemmas}
\Crefname{originallemma}{Original lemma}{Original lemmas}
\makeatletter
\newlabel{lemma-flat-local}{{03MM}{source}{}{sourcepointer-2}{}}
\newlabel{lemma-flat-local@cref}{{[originallemma][0][]03MM}{[1][0][]source}{}{}{}}
\makeatother
\crefname{originallemma}{original lemma}{original lemmas}
\Crefname{originallemma}{Original lemma}{Original lemmas}
\makeatletter
\newlabel{lemma-universally-open-local}{{03Z5}{source}{}{sourcepointer-3}{}}
\newlabel{lemma-universally-open-local@cref}{{[originallemma][0][]03Z5}{[1][0][]source}{}{}{}}
\makeatother
\crefname{originallemma}{original lemma}{original lemmas}
\Crefname{originallemma}{Original lemma}{Original lemmas}
\makeatletter
\newlabel{morphisms-lemma-fppf-open}{{01UA}{source}{}{sourcepointer-4}{}}
\newlabel{morphisms-lemma-fppf-open@cref}{{[originallemma][0][]01UA}{[1][0][]source}{}{}{}}
\makeatother
\crefname{originallemma}{original lemma}{original lemmas}
\Crefname{originallemma}{Original lemma}{Original lemmas}
\makeatletter
\newlabel{spaces-lemma-lift-morphism-presentations}{{02X1}{source}{}{sourcepointer-5}{}}
\newlabel{spaces-lemma-lift-morphism-presentations@cref}{{[originallemma][0][]02X1}{[1][0][]source}{}{}{}}
\makeatother
\title{A flat morphism locally of finite presentation of algebraic spaces is universally open}
\author{The Stacks Project Authors}
\date{stacks-project:spaces-morphisms.tex:local-snapshot-20261012; GNU FDL 1.2 or later}
\begin{document}
\maketitle
\setcounter{subsection}{1}
\noindent\textit{Editorial note.} All delivered mathematical text is the author's own, copied from the pinned Stacks Project chapter file spaces-morphisms.tex; no statement or proof step was rewritten, reordered or removed.\par\medskip
\noindent\textit{Editorial note.} The statement is one sentence with no numbered parts; the proof is the standard presentation-and-transfer argument and is complete.\par\medskip
\noindent\textit{Editorial note.} The five references of the counted spans are recorded as non-load-bearing navigation pointers; the argument that consumes them stays inside the counted proof span.\par\medskip
\noindent\textit{Editorial note.} Preamble: the three native excerpt spans of sources/stacks-spaces-morphisms/preamble.tex (lines 8-13, 14-18 and 154-217). The chapter file itself carries no preamble text.\par\medskip
\noindent\textit{Allowed standard prerequisites (editorial).}\begin{itemize}
\item Lifting a morphism of algebraic spaces to a diagram of schemes. Spaces, Lemma spaces-lemma-lift-morphism-presentations, used at proof lines 6102-6111. Produces U, V and the morphism alpha.
\item Flatness and local finite presentation are local on the source and target. Lemmas lemma-flat-local and lemma-finite-presentation-local, used at proof lines 6112-6114. Make alpha flat and locally of finite presentation.
\item The scheme case: flat plus locally of finite presentation is universally open. Morphisms, Lemma morphisms-lemma-fppf-open, used at proof lines 6115-6117. Gives universal openness of alpha.
\item The local criterion for universal openness of a morphism of algebraic spaces. Lemma lemma-universally-open-local, used at proof lines 6118-6119. Transfers universal openness from alpha to f.
\end{itemize}\medskip
\ifdefined\hypertarget\hypertarget{sourcepointer-1}{}\fi
\noindent\textit{Source reference.} Original lemma 0410: focused-seeds/AGP/sources/stacks-snapshot-20261011/tags.txt, source lines 5275--5275. Navigation pointer into the Stacks Project; the standard fact it supplies is applied inside the counted proof span and the argument resting on it stays there.\par\medskip
\ifdefined\hypertarget\hypertarget{sourcepointer-2}{}\fi
\noindent\textit{Source reference.} Original lemma 03MM: focused-seeds/AGP/sources/stacks-snapshot-20261011/tags.txt, source lines 4754--4754. Navigation pointer into the Stacks Project; the standard fact it supplies is applied inside the counted proof span and the argument resting on it stays there.\par\medskip
\ifdefined\hypertarget\hypertarget{sourcepointer-3}{}\fi
\noindent\textit{Source reference.} Original lemma 03Z5: focused-seeds/AGP/sources/stacks-snapshot-20261011/tags.txt, source lines 5208--5208. Navigation pointer into the Stacks Project; the standard fact it supplies is applied inside the counted proof span and the argument resting on it stays there.\par\medskip
\ifdefined\hypertarget\hypertarget{sourcepointer-4}{}\fi
\noindent\textit{Source reference.} Original lemma 01UA: focused-seeds/AGP/sources/stacks-snapshot-20261011/tags.txt, source lines 2403--2403. Navigation pointer into the Stacks Project; the standard fact it supplies is applied inside the counted proof span and the argument resting on it stays there.\par\medskip
\ifdefined\hypertarget\hypertarget{sourcepointer-5}{}\fi
\noindent\textit{Source reference.} Original lemma 02X1: focused-seeds/AGP/sources/stacks-snapshot-20261011/tags.txt, source lines 3824--3824. Navigation pointer into the Stacks Project; the standard fact it supplies is applied inside the counted proof span and the argument resting on it stays there.\par\medskip
\begin{lemma}
\label{lemma-fppf-open}
A flat morphism locally of finite presentation is universally open.
\end{lemma}

\begin{proof}
Let $f : X \to Y$ be a flat morphism locally of finite presentation
of algebraic spaces over $S$. Choose a diagram
$$
\xymatrix{
U \ar[r]_\alpha \ar[d] & V \ar[d] \\
X \ar[r] & Y
}
$$
where $U$ and $V$ are schemes and the vertical arrows are surjective and
\'etale, see
Spaces, Lemma \ref{spaces-lemma-lift-morphism-presentations}.
By
Lemmas \ref{lemma-flat-local} and \ref{lemma-finite-presentation-local}
the morphism $\alpha$ is flat and locally of finite presentation.
Hence by
Morphisms, Lemma \ref{morphisms-lemma-fppf-open}
we see that $\alpha$ is universally open.
Hence $X \to Y$ is universally open according to
Lemma \ref{lemma-universally-open-local}.
\end{proof}

\end{document}
